\documentclass{article}
\usepackage{amsmath}
\usepackage{titlesec}
\usepackage{graphicx}
\usepackage[margin=1in]{geometry}
\usepackage{hyperref}
\usepackage{amssymb}

% Title, date, and author
\title{Project 1}
\author{Your Name, Collaborator's Name}
\date{\today}

\titleformat{\section}
  {\normalfont\normalsize\bfseries} % Format: font style, size, and weight
  {\thesection}{1em} % Label format and spacing
  {}
  \renewcommand{\thesubsection}{\thesection.\alph{subsection}}

\titleformat{\subsection}
  {\normalfont\small\bfseries} % Format: font style, size, and weight
  {\thesubsection}{1em} % Label format and spacing
  {}
\titleformat{\subsubsection}
  {\normalfont\small\bfseries} % Format: font style, size, and weight
  {\thesubsubsection}{1em} % Label format and spacing
  {}

\begin{document}
\begin{titlepage}
    \centering
    \vspace*{1in}
    
    {\Huge\bfseries Project 1\par}
    \vspace{1.5cm}
    {\Large \today\par}
    \vspace{1.5cm}
    {\Large\itshape Antonio Pampalone 23586519 \\ Giuseppe Pisante 23610012\\ Martina Raffaelli 23616907 \par}
    
    \vfill
    \includegraphics[width=0.3\textwidth]{FAU-Logo.png}\par\vspace{1cm} % Adjust the width as needed
   
\end{titlepage}

\newpage
\small

\section*{\Large Task 2.0:}
The momentum equation is a parabolic equation since, if we compute the $\Delta = B^2 - 4AC$ from the general form of the partial differential equations: 
\[
\]
\begin{equation}
  A \frac{\partial^2 u}{\partial x^2} + B \frac{\partial^2 u}{\partial x \partial y} + C \frac{\partial^2 u}{\partial y^2} + D \frac{\partial u}{\partial x} + E \frac{\partial u}{\partial y} + F u + G = 0
\end{equation}

we get $\Delta = 0$, which means that the equation is parabolic.

We now define the additional boundary conditions:
\begin{itemize}
  \item $u(0,y) = 1$
  \item $u(1,y) = 1$
  \item $v(0,y) = 0$
  \item $v(1,y) = 0$
  \item $v(x,\infty) = 0$
\end{itemize}

\section*{\Large Task 2.1:}
The discretization applied to the x-momentum equation takes into account the parabolic nature of the equation. The discretization is performed using the central difference scheme for the y-direction, to capture the elliptic character of the diffusion process, and the backward difference scheme for the x-direction to better capture the convective term. The discretized equation is as follows:
\[
  u_{i,j}\frac{u_{i,j} - u_{i,j-1}}{\Delta x}  + v_{i,j} \frac{u_{i+1,j} - u_{i-1,j}}{2 \Delta y} = \frac{1}{Re} \frac{u_{i+1,j} - 2u_{i,j} + u_{i-1,j}}{\Delta y^2}
\]

Such equation is a steady-state equation and thus it doesn't require any time-stepping algorithm. 
For this reason, the costraints for convergence are only applied on the spatial discretization.
\\
The discretization of the first derivative in the \(x\)-direction using the Backward Difference Scheme (BDS) is given by:

\[
\frac{\partial u}{\partial x} \approx \frac{u_{i,j} - u_{i,j-1}}{\Delta x}
\]

Using Taylor series expansions for \( u(x) \) around \( x_i \), we have:

\[
u_{i,j-1} = u(x_i - \Delta x) = u(x_i) - \Delta x \frac{\partial u}{\partial x} + \frac{\Delta x^2}{2} \frac{\partial^2 u}{\partial x^2} + O(\Delta x^3)
\]

Therefore, the approximation for the derivative becomes:

\[
\frac{u_{i,j} - u_{i,j-1}}{\Delta x} = \frac{\partial u}{\partial x} + \frac{\Delta x}{2} \frac{\partial^2 u}{\partial x^2} + O(\Delta x^2)
\]

The truncation error for the Backward Difference Scheme is:

\[
T_{\text{BDS}} = -\frac{\Delta x}{2} \frac{\partial^2 u}{\partial x^2} + O(\Delta x^2)
\]

The discretization of the second derivative in the \(y\)-direction using the Central Difference Scheme (CDS) is given by:

\[
\frac{\partial^2 u}{\partial y^2} \approx \frac{u_{i+1,j} - 2u_{i,j} + u_{i+1,j}}{\Delta y^2}
\]

Using Taylor series expansions for \( u(y) \) around \( y_j \), we have:

\[
u_{i+1,j} = u(y_j + \Delta y) = u(y_j) + \Delta y \frac{\partial u}{\partial y} + \frac{\Delta y^2}{2} \frac{\partial^2 u}{\partial y^2} + \frac{\Delta y^3}{6} \frac{\partial^3 u}{\partial y^3} + O(\Delta y^4)
\]

\[
u_{i-1,j} = u(y_j - \Delta y) = u(y_j) - \Delta y \frac{\partial u}{\partial y} + \frac{\Delta y^2}{2} \frac{\partial^2 u}{\partial y^2} - \frac{\Delta y^3}{6} \frac{\partial^3 u}{\partial y^3} + O(\Delta y^4)
\]

Subtracting \( 2u_{i,j} \) from the sum of \( u_{i+1,j} \) and \( u_{i-1,j} \), we get:

\[
u_{i+1,j} - 2u_{i,j} + u_{i-1,j} = 2 \cdot \frac{\Delta y^2}{2} \frac{\partial^2 u}{\partial y^2} + O(\Delta y^4)
\]

Thus, the discretized second derivative becomes:

\[
\frac{u_{i+1,j} - 2u_{i,j} + u_{i-1,j}}{\Delta y^2} = \frac{\partial^2 u}{\partial y^2} + O(\Delta y^2)
\]

The leading truncation error for the Central Difference Scheme is:

\[
T_{\text{CDS}} = \frac{\Delta y^2}{6} \frac{\partial^4 u}{\partial y^4} + O(\Delta y^4)
\]
\subsection*{Summary of Truncation Errors}

Thus, the sum of the truncation errors is:
\[
T_{\text{total}} = T_{\text{BDS}} + T_{\text{CDS}} = -\frac{\Delta x}{2} \frac{\partial^2 u}{\partial x^2} + \frac{\Delta y^2}{6} \frac{\partial^4 u}{\partial y^4} + O(\Delta x^2) + O(\Delta y^4)
\]

\section*{\Large Task 2.2:}
This system is solved using a GMRES iterative solver, which allows us to solve the system for a non-symmetric A and the non linearities of the x-momentum.
The solver updates the values of \( u \) and \( v \) at each iteration, denoted as \( u^{(k+1)} \) and \( v^{(k+1)} \), by utilizing the values from the previous iteration, \( u^{(k)} \) and \( v^{(k)} \), to handle the non-linear terms effectively.
The system of partial differential equations is:
The continuity equation is discretized using the backward difference scheme (BDS) along the x-direction and the central difference scheme (CDS) along the y-direction. The discretized continuity equation is given by:
\[
\begin{cases}
  \frac{u_{i,j} - u_{i,j-1}}{\Delta x} + \frac{v_{i+1,j} - v_{i-1,j}}{2 \Delta y} = 0 \\
  u_{i,j}\frac{u_{i,j} - u_{i,j-1}}{\Delta x}  + v_{i,j} \frac{u_{i+1,j} - u_{i-1,j}}{2 \Delta y} = \frac{1}{Re} \frac{u_{i+1,j} - 2u_{i,j} + u_{i-1,j}}{\Delta y^2}
\end{cases}
\]


It is important to note that at the boundaries, we cannot use this discretization directly. Depending on the boundary we have to apply different schemes:
\begin{enumerate}
  \item Left boundary (\( x = 0 \)): we cannot apply BDS on the x-direction. We thus use FDS on such boundary:
  \[
\begin{cases}
  \frac{u_{i,j+1} - u_{i,j}}{\Delta x} + \frac{v_{i+1,j} - v_{i-1,j}}{2 \Delta y} = 0 \\
  u_{i,j}\frac{u_{i,j+1} - u_{i,j}}{\Delta x}  + v_{i,j} \frac{u_{i+1,j} - u_{i-1,j}}{2 \Delta y} = \frac{1}{Re} \frac{u_{i+1,j} - 2u_{i,j} + u_{i-1,j}}{\Delta y^2}
\end{cases}
\]
  
\item Right boundary (\( x = 1 \)): the problem remains the same, as BDS works well on the right boundary.
  
\item Bottom boundary (\( y = 0 \)): CDS cannot be applied on the y-direction. We thus use FDS on such boundary:
\[
  \begin{cases}
    \frac{u_{i,j} - u_{i,j-1}}{\Delta x} + \frac{v_{i,j} - v_{i-1,j}}{\Delta y} = 0 \\
    u_{i,j}\frac{u_{i,j} - u_{i,j-1}}{\Delta x}  + v_{i,j} \frac{u_{i,j} - u_{i-1,j}}{\Delta y} = \frac{1}{Re} \frac{u_{i-2,j} - 2u_{i-1,j} + u_{i,j}}{\Delta y^2}
  \end{cases}
\]
  
\item Top boundary (\( y \to \infty \)): CDS cannot be applied on the y-direction. We thus use BDS on such boundary:
\[
  \begin{cases}
    \frac{u_{i,j} - u_{i-1,j}}{\Delta x} + \frac{v_{i,j+1} - v_{i,j}}{\Delta y} = 0 \\
    u_{i,j}\frac{u_{i,j} - u_{i-1,j}}{\Delta x}  + v_{i,j} \frac{u_{i,j+1} - u_{i,j}}{\Delta y} = \frac{1}{Re} \frac{u_{i,j+2} - 2u_{i,j+1} + u_{i,j}}{\Delta y^2}
  \end{cases}
\]
\end{enumerate}

To solve this, we construct a linear system of equations in the form \( \mathbf{A} \mathbf{x} = \mathbf{b} \), where \( A \in \mathbb{R}^{2n^2 \times 2n^2} \), \( b \in \mathbb{R}^{2n^2} \), and \( x \in \mathbb{R}^{2n^2 \times 2n^2} \). The structure of \( \mathbf{x} \) is as follows:

\[
\mathbf{x} = \begin{bmatrix}
u_{1,1} \\
u_{1,2} \\
\vdots \\
u_{n,n} \\
v_{1,1} \\
v_{1,2} \\
\vdots \\
v_{n,n}
\end{bmatrix}
\]

The structure of the matrix \( \mathbf{A} \) is as follows:
\[
  \hspace{-2.5cm}
A = 
\begin{bmatrix}
  -\frac{1}{\Delta x} & \frac{1}{\Delta x} &  &  &  &  &  \vdots &  -\frac{1}{\Delta y} \cdots& \frac{1}{\Delta y} \cdots&  \\
 & \ddots  &  &  &  &  &\vdots &\ddots &  &  \\
 \cdots-\frac{1}{\Delta x} & \frac{1}{\Delta x} &  &  & & &\vdots  &\cdots -\frac{1}{2\Delta y} & &\frac{1}{2\Delta y} \cdots  \\
 & \vdots  &\ddots  &  &  &  &\vdots & &  &  \\
 &  &  &  &  &  &\vdots & &  \cdots -\frac{1}{\Delta y} & \frac{1}{\Delta y}  \\
 \cdots&\cdots  &\cdots  &\cdots  &\cdots  &\cdots  &\cdots &\cdots  &\cdots  &\cdots  \\
 -\frac{v_{i,j}^{(k)}}{\Delta y} + \frac{1}{Re \Delta y^2} - \frac{u_{i,j}^{(k)}}{\Delta x} & \frac{u_{i,j}^{(k)}}{\Delta x} &\cdots \frac{v_{i,j}^{(k)}}{\Delta y} + \frac{2}{Re \Delta y^2} & \frac{1}{Re \Delta y^2} & \cdots  &  &\vdots &  &  &  \\
 \cdots &  -\frac{1}{Re \Delta y^2}\cdots &-\frac{v_{i,j}^{(k)}}{\Delta y} + \frac{2}{Re \Delta y^2} \cdots & -\frac{u_{i,j}^{(k)}}{\Delta x}  & \frac{v_{i,j}^{(k)}}{\Delta y} - \frac{1}{Re \Delta y^2} + \frac{u_{i,j}^{(k)}}{\Delta x}  & \cdots &\vdots & & &  \\
 &\vdots&\vdots&\vdots&\vdots&\vdots&\vdots & &  &  \\
 &\cdots  & \cdots-\frac{v_{i,j}^{(k)}}{2\Delta y} + \frac{1}{Re \Delta y^2} \cdots &  -\frac{u_{i,j}^{(k)}}{\Delta x}& \frac{u_{i,j}^{(k)}}{\Delta x} - \frac{2}{Re \Delta y^2} &\cdots \frac{v_{i,j}^{(k)}}{2\Delta y} + \frac{1}{Re \Delta y^2} \cdots&\vdots &  &  &  \\
 &\vdots&\vdots&\vdots&\vdots&\vdots&\vdots & &  &  \\
 &\cdots   \frac{1}{Re \Delta y^2}\cdots & \frac{v_{i,j}^{(k)}}{\Delta y} + \frac{1}{Re \Delta y^2}\cdots & -\frac{u_{i,j}^{(k)}}{\Delta x} & \frac{u_{i,j}^{(k)}}{\Delta x} - \frac{1}{Re \Delta y^2} -\frac{v_{i,j}^{(k)}}{\Delta y} & \cdots\frac{1}{Re \Delta y^2} &\vdots &  &  & 
 

\end{bmatrix}
\]

In addition, the boundary conditions are applied to the vector \( \mathbf{b} \) to solve the system of equations, and are applied as follows:

\begin{enumerate}
  \item For \( u(0,y) = 1 \):
  \begin{itemize}
    \item Set \( b_{(j-1)n+1} = 1 \)
  \end{itemize}
  
  \item For \( u(1,y) = 1 \):
  \begin{itemize}
    \item Set \( b_{(j-1)n+n} = 1 \)
  \end{itemize}
  
  \item For \( v(0,y) = 0 \):
  \begin{itemize}
    \item Set \( b_{n^2+(j-1)n+1} = 0 \)
  \end{itemize}
  
  
   \item For \( v(1,y) = 0 \):
    \begin{itemize}
      \item Set \( b_{n^2+(j-1)n+n} = 0 \)
    \end{itemize}
  
  \item For \( v(x,\infty) = 0 \):
  \begin{itemize}
    \item Set \( b_{n^2+(n-1)n+i} = 0 \)
  \end{itemize}
\end{enumerate}

\section*{\Large Task 2.3:}
import numpy as np \\
import scipy.sparse.linalg as spla \\
M2 = spla.spilu(A)  \\
x = spla.gmres(A,b,M=M2)  \\

HINT per chatty: come utilizzare gmres per risolvere le non linearita del problema, come specificato sopra.

\begin{thebibliography}{9}
    \bibitem{GitHubRepo}
    \textit{CFD Repository},\\
    Available at: \url{https://github.com/GiuseppePisante/CFD.git}
    
    \bibitem{GitHubCopilot}
    \textit{GitHub Copilot},\\
    GitHub. Available at: \url{https://github.com/features/copilot}
    \end{thebibliography}

\end{document}